\documentclass[10pt]{article}

\usepackage[margin=1.8cm,top=1.5cm,bottom=1.5cm]{geometry}
\usepackage{xcolor}
\usepackage{hyperref}
\usepackage{amsmath}
\usepackage{enumitem}

\setlist[itemize]{noitemsep, topsep=1pt, parsep=0pt, partopsep=0pt, leftmargin=1.3em}

\hypersetup{
    colorlinks=true,
    linkcolor=blue!50!black,
    urlcolor=blue!50!black,
    citecolor=blue!50!black,
}

\newcommand{\code}[1]{\texttt{\small #1}}

\setlength{\parindent}{0pt}
\setlength{\parskip}{2pt}

\begin{document}
{\centering\large\textbf{Client Inference Attack (CIA): Threat Model}\\[2pt]
\small Federated Learning + Differential Privacy Research Project \\ \today\par}
\vspace{4pt}
\hrule
\vspace{6pt}

\small
The Client Inference Attack (CIA), introduced by S\'ainz-Pardo D\'iaz et al.\ (2026,
\emph{Knowledge-Based Systems} 343, 115993), asks whether a specific \emph{client} took part in
training a released global model --- not whether a single data record was used, as in a classic
Membership Inference Attack. Formally (Definition 8), given global model $W$ and target client
$x$'s shadow dataset $D_x$: $C(D_x, W) := P(s_x = 1 \mid D_x, W)$, where $s_x$ indicates $x$'s
participation. Motivating example: a curious pharmaceutical participant in a hospital-consortium
model wants to confirm a rival hospital's participation, as a stepping stone toward a subsequent
MIA/SIA. Implementation: \code{experiments/cia/}.

\vspace{3pt}
\textbf{Threat model --- adversary and goal}
\begin{itemize}
    \item \textbf{Who}: a \emph{semi-honest} (honest-but-curious) client that legitimately
    participates in training --- not an external party, and not the server, which is assumed
    \textbf{trusted}. The attack targets the aggregated model once released back to clients.
    \item \textbf{Goal}: infer whether a particular \emph{other} (``target'') client contributed to
    the current global model. Confirming participation narrows down whose data could explain
    observed model behavior, sharpening a follow-on MIA or SIA against that client.
    \item \textbf{Variant}: the paper's main analysis, and this repository's implementation, is the
    \emph{passive} CIA --- the attacker only observes released models, never deviates from the
    protocol. \emph{Active} (crafted/scaled updates) and \emph{adaptive} (multi-round) CIA are
    flagged as stronger future-work variants; this repo's multi-round attack extends toward the
    adaptive case empirically.
\end{itemize}

\textbf{Adversary's knowledge}
\begin{itemize}
    \item Receives the global model $W$ each round exactly as any legitimate client does --- no
    privileged access --- plus whatever aggregated per-round convergence metrics the server
    broadcasts to all clients.
    \item Central assumption: holds a shadow dataset $D_x$ approximating target $x$'s data
    distribution, either disjoint (same distribution) or overlapping (genuine target data points).
    Realistic sources: (1) public datasets partially reflecting the target's distribution,
    (2) historical data from similar sources, (3) small subsets of leaked target data.
    \item Knowledge strength is a spectrum: a fully representative shadow dataset gives an upper
    bound on attack success; a non-representative one still leaks partial signal at reduced
    confidence. Client population is assumed static across rounds (cross-silo setting).
\end{itemize}

\textbf{Adversary capabilities}
\begin{itemize}
    \item Can evaluate the released model's loss (or other metric) on its own shadow dataset each
    round it holds a checkpoint for, comparing it against the model's loss on other available data
    to estimate the target's influence on training.
    \item Cannot tamper with its own submitted updates to bias the aggregate toward revealing more
    about the target (that is the out-of-scope active-CIA variant), and cannot see other clients'
    individual updates or raw data --- only the server-released aggregate is visible.
    \item Acts on a single round's released model in the passive/first-round setting this repo's
    \code{experiments/cia} first-round attack reproduces; the repository's separate multi-round
    attack additionally exploits the trajectory of released models across rounds.
\end{itemize}

\end{document}
